\documentclass[10pt,a4paper]{article}

\usepackage[T1]{fontenc}
\usepackage[utf8]{inputenc}
\usepackage{lmodern}
\usepackage{amsmath,amssymb,amsthm}
\usepackage{booktabs,array}
\usepackage{geometry}
\usepackage{xcolor}
\usepackage{enumitem}
\usepackage{microtype}
\usepackage{hyperref}
\usepackage{fancyhdr}

\geometry{margin=1in}
\hypersetup{
  colorlinks=true,
  linkcolor=blue!55!black,
  citecolor=blue!55!black,
  urlcolor=blue!55!black,
  pdftitle={A Split-Second Forgery: A Zero-Query Attack on Random-Split St-Gen},
  pdfauthor={ePrint Signature Audit}
}
\setlist{nosep,leftmargin=*}
\setlength{\parindent}{1em}
\setlength{\parskip}{0.25em}
\pagestyle{fancy}
\fancyhf{}
\fancyhead[L]{A Split-Second Forgery}
\fancyhead[R]{9 October 2026}
\fancyfoot[C]{\thepage}
\renewcommand{\headrulewidth}{0.3pt}

\newtheorem{theorem}{Theorem}
\newtheorem{proposition}{Proposition}
\newcommand{\F}{\mathbb{F}}
\newcommand{\VES}{\mathcal{V}}

\title{\bfseries A Split-Second Forgery: A Zero-Query Attack on Random-Split St-Gen}
\author{ePrint Signature Audit\\
\small AI-assisted cryptanalysis report}
\date{9 October 2026}

\begin{document}
\maketitle

\noindent\fcolorbox{black!50}{black!4}{\parbox{0.94\textwidth}{
\textbf{AI disclosure and review status.}
OpenAI Codex generated and selected this paper's title and intentionally chose
``A Split-Second Forgery'' as wordplay on the scheme's random splitting and
the millisecond attack.  OpenAI Codex performed the scheme reconstruction,
proof development, implementation, experiments, and technical writing, with
specialist AI agents used for primary, independent, history, scope,
adversarial, and publication review.  The primary reproduction, independent
reproduction, exact-history, intended-scope, amended adversarial, and
independent publication-review gates all passed.  The result is classified
\texttt{VERIFIED\_ATTACK}.  No independent human verification was performed or
required.}}

\begin{abstract}
We describe a public-key-only, zero-query existential forgery against the
literal raw-message interface in Algorithms~5--6 of Gligoroski and
Samardjiska's random-split St-Gen signature proposal, IACR ePrint 2016/391.
The signed value specified by the paper is a tuple
$z=(z_1,z_2)\in(\F_2^n)^2$.  Verification computes
$e_i'=\sigma G_{i,\mathrm{pub}}+z_i$ and accepts when every block pair belongs
to the public \texttt{ValidErrorSplits} relation.  An attacker chooses any
$\sigma\in\F_2^k$ and public blockwise-valid $e_1,e_2$, then sets
$z_i=\sigma G_{i,\mathrm{pub}}+e_i$.  Characteristic-two cancellation makes
the verifier recover exactly the selected errors.  The result is a fresh
attacker-chosen raw-$z$ forgery with no secret input and no signing query.

Two independently structured standard-library implementations reproduce the
identity at both printed full rows: $(n,k,\ell,s)=(2142,465,3,2)$ and
$(2244,508,4,2)$.  The sealed primary run accepted zero- and nonzero-$\sigma$
variants for each row, covering all 2,550 positive blocks.  Its slowest
attack-only measurement was 1.867 milliseconds; the full run, including key
generation and controls, took 3.21 seconds and 20,992~KiB peak resident
memory.  Invalid-block, changed-signature, and representation controls reject.
Our claim is confined to the raw $z$ interface printed in Algorithms~5--6.
We make no claim for arbitrary documents or an unspecified external
document-to-$z$ hash transform, and no novelty claim for verification-first
existential forgery or linear cancellation.  The exact-history, intended-scope,
amended adversarial, and independent publication reviews passed, and the result
is classified \texttt{VERIFIED\_ATTACK}.
\end{abstract}

\section{Introduction}

Gligoroski and Samardjiska propose a code-based signature scheme whose public
key consists of two randomly split and independently permuted generator
matrices~\cite{stgen}.  The split is designed to frustrate structural and
information-set-decoding attacks considered for earlier St-Gen schemes.  The
signature verifier, however, exposes a simpler issue at the level of its
input/output relation.

Algorithm~5 receives a raw tuple $z=(z_1,z_2)$ ``to be signed.''  Algorithm~6
receives the same tuple and a signature vector $\sigma$.  Its only algebraic
operation is to recover two error vectors by adding public codewords to the
two message components.  Because the relation checked afterward is public,
an attacker can choose the checked errors first and solve the displayed
verification equations for $z$.  No inversion or decoding problem remains.

The attack has the following properties:
\begin{itemize}
  \item it applies to the exact pinned ePrint target and both complete
  parameter rows printed in Section~4;
  \item it invokes the ordinary Algorithm~6 verifier and derives the complete
  public valid-split sets from Algorithm~3;
  \item it uses no private value and makes zero signing queries;
  \item it succeeds deterministically for every public key; and
  \item it forges a fresh attacker-chosen raw $z$ tuple, the value that the
  paper itself supplies to Algorithms~5--6.
\end{itemize}

\paragraph{Classical technique attribution.}
Choosing a candidate signature first and using a public verification relation
to obtain some matching message is a classical route to existential forgery.
The standard hierarchy defines existential forgery as producing at least one
accepted message/signature pair~\cite[\S11.2.4]{hac}; the modern adaptive
chosen-message security formulation was established by Goldwasser, Micali,
and Rivest~\cite{gmr}.  We make no novelty claim for verification-first
existential forgery, linear cancellation, or the general technique.  Our
contribution is its concrete application to the intended raw-$z$ verifier in
ePrint 2016/391.  No prior public attack against the random-split St-Gen
signature scheme was located in the searches performed as of 2026-10-09.  This
is a dated search result, not an absolute priority claim.  Moody and Perlner's
earlier practical cryptanalysis applies to the unsplit predecessor and is
credited separately~\cite{moodyperlner}.

\section{Target and claim boundary}

\subsection{Pinned target}

The target is \emph{A Digital Signature Scheme Based on Random Split of
St-Gen Codes} by Danilo Gligoroski and Simona Samardjiska, IACR ePrint
2016/391~\cite{stgen}.  The exact PDF used by the audit has SHA-256 digest
\begin{center}
\small\ttfamily
b2bdf97bcef5dfec9928749d0d3fad0a97956494fc81989c52cd4472dec5755e.
\end{center}
The extracted text has SHA-256 digest
\begin{center}
\small\ttfamily
69fafe4ab6f7bb2fad7dfd6687f66dd19c5d86e738163ac5bdf75cb7fbd06d6f.
\end{center}
The reconstruction is bound to this target and to the audit's sealed primary
and independent evidence.

\subsection{Intended message space}

The unsplit predecessor, \emph{McEliece in the world of Escher}, explicitly
describes its signing input as
\texttt{Syndrome = Hash(Doc)} and contrasts that operation with the CFS loop
\texttt{Hash(Doc, Counter)}~\cite{escher}.  Its later formal algorithms use a
raw vector $z$.  The predecessor therefore proves that this lineage knew how
to state document hashing when intended.  It does not, by itself, make that
transform a mandatory unstated part of the 2016 target.

Target-specific implementation evidence resolves the remaining ambiguity.
Samy Saad Samy Shehata's 2017 NTNU master's thesis, \emph{Post Quantum
Cryptography with random split of St-Gen codes}, was supervised by Danilo
Gligoroski, a coauthor of the target~\cite{shehata}.  The thesis states that it
provides practical C implementations using the exact concrete parameter sets
from the random-split encryption and signature papers.  Its Algorithms~19 and
20 call $z_1,\ldots,z_s$ the messages, place those raw binary vectors at the
complete Sign/Verify boundary, and verify through
\[
 e_i=\sigma G_{i,\mathrm{pub}}+z_i.
\]
The appended interface passes $z$ as a byte-array argument to the signature
and verification routines, and the experimental correctness criterion is
signing a binary vector followed by successful verification.  No
document-to-$z$ transform is inserted at this interface.

The thesis is corroborating scope evidence rather than a replacement target.
It is implementation-oriented, exact-parameter work under a target author's
supervision, and it shows that the raw-vector boundary was used deliberately
after publication.  The independent scope arbitration therefore treats the
raw binary-vector space as the intended mathematical signature interface.  A
preimage-resistant external wrapper cannot be imported from the predecessor as
a mandatory hidden step.

\subsection{Public error-split relation}

Let $E_\ell\subseteq\F_2^\ell$ be the error set and let $s=2$ be the number
of matrix splits.  Definition~3 calls $(a,b)\in(\F_2^\ell)^2$ a valid error
split when
\[
  \pi_1(a)+\pi_2(b)\in E_\ell
\]
for every pair of coordinate permutations $\pi_1,\pi_2$.  Algorithm~3
enumerates all such pairs once, and the paper says the resulting
\texttt{ValidErrorSplits} set may be public.  Write that set as $\VES$.

The two printed rows use error sets consisting of all $\ell$-bit words of
Hamming weight one or two.  Exhaustive implementation of Algorithm~3 gives
the counts stated by the paper:
\begin{center}
\begin{tabular}{lrrrrr}
\toprule
Row & $n$ & $k$ & $\ell$ & blocks $n/\ell$ & $|\VES|$ \\
\midrule
PS1 & 2142 & 465 & 3 & 714 & 24 \\
PS2 & 2244 & 508 & 4 & 561 & 40 \\
\bottomrule
\end{tabular}
\end{center}

\subsection{Algorithms 5 and 6}

For $s=2$, Algorithm~5 takes
\[
 z=(z_1,z_2)\in(\F_2^n)^2
\]
as the value to be signed and returns $\sigma\in\F_2^k$.  It intends the
resulting errors
\[
 e_i=\sigma G_{i,\mathrm{pub}}+z_i
\]
to have every corresponding block pair in $\VES$.

Algorithm~6 receives $(z,\sigma)$ and the two public matrices.  It computes
\begin{equation}
 e_i'=\sigma G_{i,\mathrm{pub}}+z_i,
 \qquad i\in\{1,2\},
 \label{eq:verify}
\end{equation}
and accepts precisely when every $\ell$-bit block pair of $(e_1',e_2')$
belongs to $\VES$.

The target does not define a document-to-$z$ encoding, a preimage-resistant
hash-to-$z$ transform, or a byte-level external-message format.  Our result is
therefore limited to the intended raw tuple accepted by the printed algorithms
and corroborated by the thesis.  We do not claim a forgery for arbitrary
documents, a preimage under a separately defined wrapper, the ePrint 2014/360
predecessor, or any software implementation.

\section{Zero-query forgery}

\begin{theorem}
For every public key generated by Algorithm~2, every
$\sigma\in\F_2^k$, and every pair of vectors $(e_1,e_2)$ whose corresponding
$\ell$-bit blocks all lie in $\VES$, the tuple
\begin{equation}
 z_i=\sigma G_{i,\mathrm{pub}}+e_i,
 \qquad i\in\{1,2\},
 \label{eq:forge}
\end{equation}
is accepted with signature $\sigma$ by Algorithm~6.
\end{theorem}

\begin{proof}
Substituting Equation~\eqref{eq:forge} into the verifier's
Equation~\eqref{eq:verify} gives
\[
 e_i'=\sigma G_{i,\mathrm{pub}}+
      \sigma G_{i,\mathrm{pub}}+e_i=e_i,
\]
because each field element has characteristic two.  Every checked block pair
was selected from $\VES$, so Algorithm~6 accepts.
\end{proof}

\paragraph{Attack algorithm.}
The attacker enumerates $\VES$ from public parameters, fills all $n/\ell$
blocks of $(e_1,e_2)$ with valid pairs, chooses $\sigma$, computes the two
public matrix products, and outputs Equation~\eqref{eq:forge}.

The choice $\sigma=0$ reduces the attack to $z_i=e_i$ and shows that the
acceptance identity is independent of the key.  A nonzero $\sigma$ variant
exercises the complete public-matrix path and demonstrates that excluding the
zero signature would not help.  The result holds for every chosen $\sigma$.

\paragraph{Freshness and attack class.}
The raw tuple $z$ is computed for the attack and never submitted to a signing
oracle; there is no signing query at all.  Thus the output is a fresh
message/signature pair for the literal message space.  Because the attacker
chooses the signature and obtains a matching message rather than targeting a
preselected external document, the precise class is a zero-query existential
forgery.  That class is sufficient to violate existential unforgeability of
the raw-message construction.

\section{Primary implementation}

The primary standard-library implementation represents binary row vectors as
Python integers.  It reconstructs all relevant paper algorithms rather than
hard-coding an accepted transcript:
\begin{enumerate}
  \item It constructs the Equation~(1) staircase generator $G$ for the full
  PS1 and PS2 dimensions.
  \item It samples $G_1$ as a full binary matrix and defines $G_2=G+G_1$.
  \item It applies an invertible scrambler, a common outer block permutation,
  and independent inner permutations to form both public matrices.
  \item It exhaustively derives $\VES$ from Algorithm~3's universal
  permutation condition.
  \item A public-only function produces zero- and nonzero-$\sigma$ forgeries.
  \item The ordinary Algorithm~6 function checks every block.
  \item Invalid-block and strict representation controls test rejection.
\end{enumerate}

The deterministic seeds make the public-package output reproducible and are
not used by the attack.  The forged-message digests and public-key digests are
stable.  The private staircase matrix, split, scrambler, and permutations are
held separately and are never passed to the forger.

\subsection{Primary results}

The sealed primary run produced the following measurements:
\begin{center}
\small
\begin{tabular}{llrrrrl}
\toprule
Row & variant & blocks & $|\VES|$ & attack & verify & result \\
\midrule
PS1 & $\sigma=0$ & 714 & 24 & 1.867 ms & 0.420 ms & accept \\
PS1 & nonzero $\sigma$ & 714 & 24 & 0.469 ms & 0.504 ms & accept \\
PS2 & $\sigma=0$ & 561 & 40 & 0.423 ms & 0.460 ms & accept \\
PS2 & nonzero $\sigma$ & 561 & 40 & 0.501 ms & 0.523 ms & accept \\
\bottomrule
\end{tabular}
\end{center}

All four variants cover the complete valid-split set for their row.  The
ordinary verifier checked all 2,550 positive blocks.  Four one-invalid-block
controls and sixteen format controls rejected.  The complete run, including
both key generations and all controls, took 3.21 seconds wall-clock and used
20,992~KiB peak resident memory.  The maximum attack-only measurement was
1.867 milliseconds.

\section{Independent reproduction}

The sealed independent reproduction was written separately from the primary
implementation.  It used OS-random full-row keys and a dense random full-rank
scrambler.  Key generation, the attacker, and verification ran as separate
processes.  The key generator persisted no private material, and the attacker
process received only a serialized public key and an output path; it exposed
no signing-oracle interface.

For one fresh key at each full row, the independent verifier accepted both a
zero-$\sigma$ and a nonzero-$\sigma$ forgery.  It rejected an invalid block, a
changed raw message, a changed signature, and four representation controls at
each row.  Peak resident memory was 20,480~KiB, and each of the six stages took
at most 0.66 seconds.  The sealed manifest has SHA-256 digest
\begin{center}
\small\ttfamily
c0e2c712fccb11803285f2f0bf03c16cee226b7b15f91d0eaab9ece71fae064d.
\end{center}

The package also contains a deterministic, separately structured companion
reproducer.  It uses a dense full-rank scrambler and independent matrix,
permutation, enumeration, attack, and verification routines.  It repeats both
variants and the controls at both rows, yielding stable reference output.

\section{Controls and interpretation}

\subsection{Invalid-block controls}

At both rows, $(0,0)$ is absent from $\VES$.  The controls replace one block
of $z$ so that Algorithm~6 recovers exactly this pair.  The primary zero
variants fail at the first block; the nonzero variants place the bad block at
the end and therefore exercise the complete rejection scan.  The independent
implementation also changes a valid signature by one bit, selected by an
exhaustive search, while keeping $z$ fixed.  Verification rejects.

\subsection{Representation controls}

The mathematical algorithms declare fixed-dimensional vector spaces.  The
implementations enforce those dimensions explicitly.  Wrong message arity,
wrong component width, wrong signature width, non-hexadecimal encodings in
the sealed independent run, and values outside the declared vector spaces all
reject before the algebraic predicate.  These tests rule out an accidental
always-accepting verifier and separate the identity from encoding behavior.

\subsection{Separate signer-correctness observation}

The primary implementation also records a correctness gap that the forgery
does not use.  Algorithm~5's stated decoder contract ensures only that an
aggregate decoded error belongs to $E_\ell$.  Algorithm~6 demands the stronger
condition that each component tuple belongs to $\VES$, which quantifies over
every independent pair of inner permutations.

Take two distinct unit vectors $a,b\in\F_2^\ell$.  Their sum has Hamming
weight two and belongs to the printed $E_\ell$.  Independent permutations can
align the two units, however, making their sum zero, which is outside
$E_\ell$.  Thus $(a,b)\notin\VES$.  The primary run constructs key-consistent
full-row transcripts that meet the stated aggregate decoder condition at all
714 PS1 blocks or all 561 PS2 blocks, yet Algorithm~6 rejects every block.
This is separate from the attack, which directly constructs accepting
transcripts.

\section{Limitations and repair direction}

The result is exact for the mathematical raw-$z$ interface in the pinned
paper.  It does not establish an attack on an absent external transform from
meaningful documents to $z$.  In particular, if a deployment independently
specified a preimage-resistant map from documents into the complete $z$
space, the attacker would additionally have to find a document mapping to the
constructed tuple.  That is a different problem and no such transform appears
in the target.

The experiment does not estimate a success probability.  Once public valid
errors are selected, acceptance is an algebraic identity for every key.  The
fixed seeds in the package affect only reproducibility of test keys.

Banning $\sigma=0$ is insufficient because Equation~\eqref{eq:forge} works
for every nonzero $\sigma$.  A repair must define the signed message domain
and prevent an adversary from solving backward from a freely chosen signature
to an accepted signed value.  An explicit, analyzed preimage-resistant
encoding from external messages into $z$ could block this particular route,
but it would be a new transform requiring its own correctness and security
proof.  More generally, the verification relation should bind the accepted
object to the signing key under a standard unforgeability reduction rather
than expose a publicly solvable affine equation.

\section{Conclusion}

For both printed full parameter rows, a zero-query adversary can choose public
valid error blocks and any signature vector, then solve the verifier equation
for a fresh raw $z$ tuple.  Algorithm~6 cancels the duplicated public
codeword and accepts deterministically.  Two independently structured
implementations reproduce the attack, and the measured construction takes at
most 1.867 milliseconds in the sealed primary run.  The exact claim is a
public-key-only existential forgery against the raw tuple specified by
Algorithms~5--6.  It does not extend to arbitrary documents or an unspecified
external hash wrapper.  Primary, independent, history, intended-scope,
amended adversarial, and independent publication reviews all passed.  The
result is classified \texttt{VERIFIED\_ATTACK}.

\section*{Acknowledgements}
The computations were performed on the Norwegian Research and Education Cloud (NREC), using resources provided by the University of Bergen and the University of Oslo.

\begin{thebibliography}{9}

\bibitem{stgen}
D.~Gligoroski and S.~Samardjiska.
\newblock A Digital Signature Scheme Based on Random Split of St-Gen Codes.
\newblock IACR Cryptology ePrint Archive, Report 2016/391, 2016.
\newblock \url{https://eprint.iacr.org/2016/391}.

\bibitem{escher}
D.~Gligoroski, S.~Samardjiska, H.~Jacobsen, and S.~Bezzateev.
\newblock McEliece in the world of Escher.
\newblock IACR Cryptology ePrint Archive, Report 2014/360, 2014.
\newblock \url{https://eprint.iacr.org/2014/360}.

\bibitem{shehata}
S.~S.~S.~Shehata.
\newblock \emph{Post Quantum Cryptography with random split of St-Gen codes}.
\newblock Master's thesis, Norwegian University of Science and Technology,
June 2017.  Supervised by Danilo Gligoroski.
\newblock \url{https://hdl.handle.net/11250/2450584}.

\bibitem{moodyperlner}
D.~Moody and R.~Perlner.
\newblock Vulnerabilities of ``McEliece in the World of Escher.''
\newblock In \emph{Post-Quantum Cryptography 2016}, LNCS 9606, pp.~104--117.
Springer, 2016.
\newblock \url{https://doi.org/10.1007/978-3-319-29360-8_8}.

\bibitem{gmr}
S.~Goldwasser, S.~Micali, and R.~L.~Rivest.
\newblock A digital signature scheme secure against adaptive chosen-message
attacks.
\newblock \emph{SIAM Journal on Computing}, 17(2):281--308, 1988.
\newblock \url{https://doi.org/10.1137/0217017}.

\bibitem{hac}
A.~J.~Menezes, P.~C. van Oorschot, and S.~A. Vanstone.
\newblock \emph{Handbook of Applied Cryptography}, Chapter~11.
\newblock CRC Press, 1996.
\newblock \url{https://cacr.uwaterloo.ca/hac/about/chap11.pdf}.

\end{thebibliography}

\end{document}
